\documentclass[12pt,a4paper]{article}

\usepackage[utf8]{inputenc}
\usepackage[T1]{fontenc}
\usepackage[margin=2.5cm]{geometry}
\usepackage{hyperref}
\usepackage{microtype}
\usepackage{parskip}
\usepackage{lmodern}

\hypersetup{colorlinks=true, urlcolor=blue!60!black}

\pagestyle{empty}

\begin{document}

\begin{flushright}
Kundai Farai Sachikonye \\
Auf der Lichtung 19, 82178 Puchheim, Germany \\
\href{https://github.com/fullscreen-triangle}{github.com/fullscreen-triangle} \\
\texttt{kundai.sachikonye@bitspark.com} \\
(+49) 1626386344 \\[6pt]
18 May 2026
\end{flushright}

\vspace{1em}

Prof.\ Dr.\ Anna Poetsch \\
Poetsch Group --- Institut für Genetik / CECAD \\
Universität zu Köln \\
\texttt{apoetsch@uni-koeln.de} \\
Ref.\ Wiss2604-09

\vspace{1.5em}

\textbf{Re: Deep Learning Engineer --- Omics Data / Genomics (Ref.\ Wiss2604-09)}

\vspace{1em}

Dear Prof.\ Poetsch,

I am applying for the Deep Learning Engineer position in the Poetsch group
(Ref.\ Wiss2604-09). The position describes building deep learning
infrastructure on omics data for the lab and the broader research community ---
converting theoretical concepts into working code, implementing state-of-the-art
algorithms, and creating internal tools that accelerate research.
This is exactly the kind of work I have been doing independently for four years,
across genomics, transcriptomics, mass spectrometry, and immunopeptidology data.
The frameworks I have built are not academic exercises; they are deployed,
validated, and operational systems that I can bring into the lab's infrastructure.

\textbf{Deep learning for genomics and omics data.}
The \href{https://github.com/fullscreen-triangle/gospel}{Gospel} framework
implements the full genomics analysis pipeline at production scale:
$10^6$ variants per second VCF processing, RNA-seq quantification and differential
expression, Bayesian network inference for multi-omics integration, and Gaussian
mixture models for population structure.
Validated against 1000~Genomes Phase~3, gnomAD~v3.1.2, and UK~Biobank.
For the Poetsch group's work on genomic changes in aging and cancer, Gospel
provides the computational backbone for somatic variant calling,
tumour mutational burden analysis, and integrating genomic with transcriptomic
data across patient cohorts.

The \href{https://github.com/fullscreen-triangle/lavoisier}{Lavoisier} framework
demonstrates the specific DL pattern the position describes --- implementing and
evaluating a state-of-the-art algorithm from a recent academic paper and deploying
it as a research tool. The dual-pipeline design combines:
\begin{itemize}
  \item A numerical pipeline (standard MS1/MS2 processing, peak detection,
    spectral library matching) validated against NIST at 98.9\,\% accuracy.
  \item A CNN-based visual pipeline (PyTorch) that converts mass spectra
    into temporal image sequences and applies image classification
    networks --- moving from the academic literature on spectral imaging
    into a working, distributed production system (Ray, $>$100\,GB mzML support).
\end{itemize}
This is precisely the "from academic papers to working code" task the position
lists as a core responsibility. The same architectural pattern --- parallel
classical and deep learning pipelines with cross-validated accuracy metrics ---
applies directly to genomics data: a classical variant calling pipeline alongside
a CNN or Transformer-based variant pathogenicity predictor, both sharing the
same input pre-processing layer.

\textbf{Algorithm implementation and internal tooling.}
The position emphasises creating internal tools that speed up research:
automated evaluation frameworks, data annotation tools, specialized libraries.
I have built all three categories:

\textit{Automated evaluation frameworks:} Every framework I have built includes
automated validation against reference standards with quantified accuracy metrics.
The Lavoisier pipeline runs automated benchmarking against the NIST mass spectral
library on each build; Gospel validates variant call accuracy against gnomAD
population frequencies and Clinvar annotations; the Syndrome immunopeptidology
tools benchmark against the IEDB peptide database.

\textit{Data annotation tools:} The
\href{https://syndrome-seven.vercel.app/tools}{Syndrome} live tools provide
browser-based annotation for genomic sequences --- MHC binding prediction,
neoantigen identification, immune escape scoring --- requiring no local
installation. This is the kind of interactive annotation tool that lab members
with varying computational expertise can use directly.

\textit{Specialised libraries:} The Rust core of
\href{https://github.com/fullscreen-triangle/four-sided-triangle}{Four-Sided-Triangle}
is a specialised library for high-throughput Bayesian inference, fuzzy logic,
and evidence propagation, providing 15--50$\times$ speedup over Python via
PyO3 bindings. The same library architecture --- Rust performance-critical
operations exposed to Python via clean API --- is the standard approach for
building genomics DL libraries where inference throughput matters.

\textbf{Aging and cancer genomics: domain context.}
The Poetsch group's focus on genomic changes in aging and cancer is the
application domain I have studied most intensively from a computational
perspective.
For cancer: the Syndrome framework's four live immunopeptidology tools
(neoantigen identification, MHC binding, immune escape, drug resistance)
directly analyse the somatic mutations and immune interactions relevant to
cancer genomics. The Crown Prince platform derives tumour cell diagnostics
from spectral hologram representations, with AUC$\,=1.00$ on the loop-holonomy
self-consistency diagnostic.
For aging: the \href{https://rembling-trambling.vercel.app}{Olduvai}
framework analyses biological motion degradation patterns across aging
(Parkinson's, ataxia, healthy ageing), validated on 86 nights of wearable
sensor data. The D$\in[0,1]^8$ disease vector in Syndrome has an explicit
circadian regulatory oscillator dimension ($D_\mathrm{Circ}$) and ATP synthesis
coherence dimension ($D_\mathrm{ATP}$) --- the two cellular dimensions most
directly coupled to the aging process at the genomic expression level.
Connecting genomic variation data (from Gospel's variant analysis) to these
cellular phenotype dimensions is the kind of integrative analysis the position
is structured to enable.

\textbf{Equivalent experience in place of a completed doctorate.}
The position specifies a PhD in CS, Bioinformatics, or AI ``or equivalent
experience level.''
I completed doctoral research in Computational Lipidomics at TUM /
Bayerisches Forschungsinstitut (2019--2021) and have since worked full-time
as an independent researcher, producing peer-reviewed-quality manuscripts
(three submitted, affiliated with AIMe Registry and TUM Weihenstephan) and
building twelve deployed, validated computational platforms.
The depth and scope of this independent research is the equivalent experience
the position's clause refers to.

I work primarily in Python with PyTorch; Rust for performance-critical
components; Ray for distributed computing; Docker and Kubernetes for deployment.
English is native; German is conversational.
I am legally resident in Germany with full work authorisation and available
from July 2026.

\vspace{1.5em}
Yours sincerely,

\vspace{2.5em}
\textbf{Kundai Farai Sachikonye}

\end{document}
